本附录回顾正文中反复使用的若干基础材料.除特别说明外,\(\mathbb F\) 表示 \(\mathbb R\) 或 \(\mathbb C\).

\section{向量空间的张量积}

给定两个复向量空间 \(V_1\) 与 \(V_2\),张量积是一个新的向量空间 \(V_1\otimes V_2\),并配有一个双线性“乘积”映射
\(\otimes:V_1\times V_2\to V_1\otimes V_2\).若 \(V_1,V_2\) 为有限维空间,且 \(\{u_j\}\)、\(\{v_k\}\) 分别为它们的基,则 \(V_1\otimes V_2\) 为有限维空间,而 \(\{u_j\otimes v_k\}\) 构成其一组基.有限维情形当然可以直接用这一基性质定义张量积;但这样做必须另外说明构造与所选基无关.因此,我们采用下面的“泛性质”来定义张量积.

\begin{definition}[张量积的泛性质]\label{app:def:tensor}
设 \(V_1,V_2\) 是域 \(\mathbb F\) 上的向量空间.称向量空间 \(W\) 连同双线性映射 \(T:V_1\times V_2\to W\) 为 \(V_1,V_2\) 的一个张量积,如果满足下述泛性质:对任意 \(\mathbb F\) 上的向量空间 \(U\) 以及双线性映射 \(\Phi:V_1\times V_2\to U\),存在唯一的线性映射 \(\widetilde\Phi:W\to U\),使得相应的图交换.也就是说,\(\widetilde\Phi\circ T=\Phi\).
\end{definition}

\begin{proposition}[张量积的存在性与唯一性]\label{app:prop:tensor-existence}
任意两个向量空间 \(V_1,V_2\) 都存在张量积,并且张量积在“典范同构”的意义下唯一.具体而言,若 \((W_1,T_1)\) 与 \((W_2,T_2)\) 都是张量积,则存在唯一的可逆线性映射 \(\Psi:W_1\to W_2\),满足 \(T_2=\Psi\circ T_1\).
\end{proposition}

由此唯一性结果,我们可以直接谈论“这个”张量积,并把它记作 \(V_1\otimes V_2\).同时把线性映射 \(T:V_1\times V_2\to V_1\otimes V_2\) 写成 \((u,v)\mapsto u\otimes v\).在此记号下,泛性质表示为:给定从 \(V_1\times V_2\) 到向量空间 \(U\) 的任意双线性映射 \(\Phi\),存在唯一的线性映射 \(\widetilde\Phi:V_1\otimes V_2\to U\),使得
\[
\widetilde\Phi(u\otimes v)=\Phi(u,v).
\]

\begin{proposition}[有限维张量积的基]\label{app:prop:tensor-basis}
设 \(V_1,V_2\) 是有限维向量空间,\(\{u_j\}_{j=1}^{n_1}\)、\(\{v_k\}_{k=1}^{n_2}\) 分别为它们的基.则 \(V_1\otimes V_2\) 为有限维空间,所有形如 \(u_j\otimes v_k\)(\(1\le j\le n_1,\ 1\le k\le n_2\))的元素构成其一组基.特别地,
\[
\dim(V_1\otimes V_2)=(\dim V_1)(\dim V_2).
\]
\end{proposition}

需要强调的是,一般而言,\(V_1\otimes V_2\) 中并非每个元素都具有 \(u\otimes v\) 的形式,其中 \(u\in V_1,v\in V_2\).我们只能说,每个元素都能写成若干形如 \(u\otimes v\) 的元素的线性组合,而且这种分解远非典范;即使在有限维情形,它也依赖于 \(V_1,V_2\) 中基的选择.然而,张量积的泛性质告诉我们,只要在形如 \(u\otimes v\) 的元素上给出定义,就可以构造从 \(V_1\otimes V_2\) 到任意向量空间 \(U\) 的线性映射.只要 \(\Phi(u,v)\) 关于 \(u,v\) 双线性,泛性质就保证存在唯一的线性映射 \(\widetilde\Phi\),使得 \(\widetilde\Phi(u\otimes v)=\Phi(u,v)\).下面的结果是泛性质的一个典型应用.

\begin{proposition}[算子的张量积]\label{app:prop:tensor-operators}
若 \(A\in\operatorname{End}(V_1)\)、\(B\in\operatorname{End}(V_2)\),则存在唯一的线性映射 \(A\otimes B:V_1\otimes V_2\to V_1\otimes V_2\),满足
\[
(A\otimes B)(u\otimes v)=(Au)\otimes(Bv).
\]
此外,对 \(A_1,A_2\in\operatorname{End}(V_1)\) 与 \(B_1,B_2\in\operatorname{End}(V_2)\),有
\[
(A_1\otimes B_1)(A_2\otimes B_2)=(A_1A_2)\otimes(B_1B_2).
\]
\end{proposition}

要构造 \(A\otimes B\),只需在泛性质中取 \(U=V_1\otimes V_2\),并令 \(\Phi(u,v)=(Au)\otimes(Bv)\).由于 \(A,B\) 是线性的,而张量积是双线性的,\(\Phi\) 是双线性映射;于是泛性质给出的线性映射 \(\widetilde\Phi\) 就是 \(A\otimes B\).

这里定义的张量积适用于有限维和无限维向量空间.不过,这个构造纯粹是代数性的;若 \(V_1,V_2\) 带有拓扑,张量积并不考虑这些拓扑.因此,在 Hilbert 空间的背景下,我们还需要细化张量积的概念,使两个 Hilbert 空间的张量积仍为 Hilbert 空间,见附录 \ref{app:subsec:hilbert-tensor}.

\section{测度论}

本节假定读者熟悉测度论的基本概念,包括 \(\sigma\)-代数、测度、可测函数与 Lebesgue 积分.

三元组 \((X,\Omega,\mu)\)(其中 \(X\) 是集合,\(\Omega\) 是 \(X\) 的子集构成的 \(\sigma\)-代数,\(\mu\) 是 \(\Omega\) 上的非负测度)称为测度空间.

可测函数 \(\psi:X\to\mathbb C\) 称为可积的,如果 \(\int_X|\psi|\,d\mu<\infty\).由 \(X\) 的任意一族子集生成的 \(\sigma\)-代数,是包含这族子集的最小 \(\sigma\)-代数.

我们把单调收敛定理、控制收敛定理、\(L^p\) 空间以及 Fubini 定理视为标准知识.下面简要回顾几项可能不那么熟悉的内容.

若测度空间 \((X,\Omega,\mu)\) 中的 \(X\) 能写成可数个有限测度可测集的并,则称测度 \(\mu\) 是 \(\sigma\)-有限的.

\begin{definition}[绝对连续与等价测度]\label{app:def:absolute-continuity}
设 \(\mu,\nu\) 是测度空间 \((X,\Omega)\) 上的两个 \(\sigma\)-有限测度.若对每个 \(E\in\Omega\),\(\nu(E)=0\) 都推出 \(\mu(E)=0\),则称 \(\mu\) 关于 \(\nu\) 绝对连续.若两个测度彼此绝对连续,则称它们等价.
\end{definition}

\begin{theorem}[Radon--Nikodym 定理]\label{app:thm:radon-nikodym}
设 \(\mu,\nu\) 是测度空间 \((X,\Omega)\) 上的两个 \(\sigma\)-有限测度,且 \(\mu\) 关于 \(\nu\) 绝对连续.则存在 \(X\) 上的非负可测函数 \(\rho\),使得对所有 \(E\in\Omega\),
\[
\mu(E)=\int_E\rho\,d\nu.
\]
函数 \(\rho\) 称为 \(\mu\) 关于 \(\nu\) 的密度.
\end{theorem}

\begin{definition}[单调类]\label{app:def:monotone-class}
集合 \(X\) 的子集族 \(\mathcal M\) 称为单调类,如果它对可数递增并与可数递减交封闭.可数递增并是指一列满足 \(E_j\subset E_{j+1}\) 的集合的并;可数递减交类似定义.
\end{definition}

\begin{theorem}[单调类引理]\label{app:thm:monotone-class}
设 \(\mathcal M\) 是 \(X\) 的一个单调类,并且包含 \(X\) 的一个代数 \(\mathcal A\).则 \(\mathcal M\) 包含由 \(\mathcal A\) 生成的 \(\sigma\)-代数.
\end{theorem}

\begin{corollary}\label{app:cor:measure-agreement}
设 \(\mu,\nu\) 是测度空间 \((X,\Omega)\) 上的两个有限测度,并且它们在代数 \(\mathcal A\subset\Omega\) 上相等,则它们在由 \(\mathcal A\) 生成的 \(\sigma\)-代数上也相等.
\end{corollary}

\begin{note}
一般而言,两个测度取值相等的集合所组成的族既不一定是 \(\sigma\)-代数,甚至不一定是代数.
\end{note}

\begin{theorem}[紧支集连续函数的稠密性]\label{app:thm:cc-dense}
设 \(\mu\) 是局部紧、可分度量空间 \(X\) 的 Borel \(\sigma\)-代数上的测度,并且每个紧子集 \(K\) 都满足 \(\mu(K)<\infty\).那么,对所有 \(1\le p<\infty\),\(X\) 上紧支集连续函数空间在 \(L^p(X,\mu)\) 中稠密.
\end{theorem}

这里需要澄清一句话.若 \(\psi\) 是 \(X\) 上的紧支集连续函数,则由于 \(\psi\) 有界且 \(\mu\) 在紧集上有限,\(\int_X|\psi|^p\,d\mu\) 有限.因此,我们可以把 \(C_c(X)\) 中的 \(\psi\) 映到 \(L^p(X,\mu)\) 中的等价类 \([\psi]\).定理更准确的表述是:这个映射的像在 \(L^p(X,\mu)\) 中稠密.不过,\(\psi\mapsto[\psi]\) 不必是单射.例如,若 \(X\) 中有一个非空开集 \(U\) 满足 \(\mu(U)=0\),那么任何支集包含于 \(U\) 的 \(\psi\) 都满足 \([\psi]=0\).尽管如此,我们仍会略微滥用术语,称 \(C_c(X)\) 在 \(L^p(X,\mu)\) 中稠密.

\section{初等泛函分析}\label{app:sec:elementary-functional-analysis}

本节简要回顾正文中用到的初等泛函分析结果.大多数结果都可以在 Rudin 的著作 \cite{RudinReal,RudinFunctional} 中找到.

\subsection{Stone--Weierstrass 定理}

Weierstrass 定理断言,任意区间上的连续实值函数都可以由多项式一致逼近.Stone 对此作出了重要推广.设 \(X\) 是紧度量空间,记 \(C(X;\mathbb R)\)、\(C(X;\mathbb C)\) 为连续实值函数空间与连续复值函数空间.\(C(X;\mathbb F)\) 的子集 \(\mathcal A\) 称为代数,如果它对逐点加法、逐点乘法以及乘以 \(\mathbb F\) 的标量封闭.若任意两个不同的点 \(x,y\in X\) 都存在 \(f\in\mathcal A\) 使 \(f(x)\ne f(y)\),则称 \(\mathcal A\) 分离点.我们在 \(C(X;\mathbb F)\) 上使用上确界范数
\[
\|f\|_{\mathrm{sup}}:=\sup_{x\in X}|f(x)|,
\]
并且 \(C(X;\mathbb F)\) 关于距离 \(d(f,g)=\|f-g\|_{\mathrm{sup}}\) 完备.

\begin{theorem}[Stone--Weierstrass 定理,实版本]\label{app:thm:stone-real}
设 \(X\) 是紧度量空间,\(\mathcal A\) 是 \(C(X;\mathbb R)\) 中的代数.若 \(\mathcal A\) 包含常值函数并且分离点,则 \(\mathcal A\) 关于上确界范数在 \(C(X;\mathbb R)\) 中稠密.
\end{theorem}

\begin{theorem}[Stone--Weierstrass 定理,复版本]\label{app:thm:stone-complex}
设 \(X\) 是紧度量空间,\(\mathcal A\) 是 \(C(X;\mathbb C)\) 中的代数.若 \(\mathcal A\) 包含常值函数、分离点并且对复共轭封闭,则 \(\mathcal A\) 关于上确界范数在 \(C(X;\mathbb C)\) 中稠密.
\end{theorem}

复版本的 Stone--Weierstrass 定理有如下推论:若 \(K\) 是 \(\mathbb C\) 的紧子集,则 \(K\) 上任意连续复值函数都可以由 \(z\) 与 \(\bar z\) 的多项式一致逼近.

\subsection{Fourier 变换}\label{app:subsec:fourier-transform}

现在以几种不同形式描述 \(\mathbb R^n\) 上的 Fourier 变换.

\begin{definition}[Fourier 变换]\label{app:def:fourier}
对任意 \(\psi\in L^1(\mathbb R^n)\),定义其 Fourier 变换为 \(\mathbb R^n\) 上的函数
\[
\widehat\psi(k)=(2\pi)^{-n/2}\int_{\mathbb R^n}e^{-ik\cdot x}\psi(x)\,dx.
\]
\end{definition}

\begin{proposition}\label{app:prop:fourier-basic}
若 \(\psi\in L^1(\mathbb R^n)\),则其 Fourier 变换满足:\(|\widehat\psi(k)|\le(2\pi)^{-n/2}\|\psi\|_{L^1}\),\(\widehat\psi\) 连续,并且当 \(|k|\to\infty\) 时,\(\widehat\psi(k)\to0\).
\end{proposition}

前两项分别由定义和控制收敛定理立即得到.为证明第三项,先在 \(L^1(\mathbb R^n)\) 的一个稠密子空间(例如下面定义的 Schwartz 空间)上证明,再取一致极限.

\begin{definition}[Schwartz 空间]\label{app:def:schwartz}
\(\mathcal S(\mathbb R^n)\) 是所有 \(C^\infty\) 函数 \(\psi\) 的空间,它们满足:对任意由非负整数构成的 \(n\)-元组 \(j,k\),均有
\[
\lim_{x\to\pm\infty}|x^j\partial^k\psi(x)|=0.
\]
这里若 \(j=(j_1,\ldots,j_n)\),则 \(x^j=x_1^{j_1}\cdots x_n^{j_n}\),并且
\[
\partial^k=\left(\frac{\partial}{\partial x_1}\right)^{k_1}\cdots\left(\frac{\partial}{\partial x_n}\right)^{k_n}.
\]
Schwartz 空间中的元素称为 Schwartz 函数.
\end{definition}

\begin{proposition}\label{app:prop:fourier-schwartz}
若 \(\psi\in\mathcal S(\mathbb R^n)\),则 \(\widehat\psi\in\mathcal S(\mathbb R^n)\).
\end{proposition}

这个结果的证明依赖于 Fourier 变换在微分以及乘以 \(x\) 时的行为,下面的结果正是这两种行为.

\begin{proposition}\label{app:prop:fourier-derivatives}
若 \(\psi\) 是 Schwartz 函数,则
\[
\widehat{\frac{\partial\psi}{\partial x_j}}(k)=ik_j\widehat\psi(k),\tag{A.1}
\]
并且 \(\widehat\psi\) 在每一点可微,且
\[
\widehat{x_j\psi}(k)=i\frac{\partial}{\partial k_j}\widehat\psi(k).\tag{A.2}
\]
第一式由分部积分证明,第二式由在 Fourier 变换定义下对积分求导证明.
\end{proposition}

\begin{theorem}[Fourier 反演与 Plancherel 公式 I]\label{app:thm:plancherel-I}
在 \(\mathcal S(\mathbb R^n)\) 上,Fourier 变换把 Schwartz 空间映到自身且满射,并且对任意 \(\psi\in\mathcal S(\mathbb R^n)\),
\[
\psi(x)=(2\pi)^{-n/2}\int_{\mathbb R^n}e^{ik\cdot x}\widehat\psi(k)\,dk,
\]
以及
\[
\int_{\mathbb R^n}|\psi(x)|^2\,dx=\int_{\mathbb R^n}|\widehat\psi(k)|^2\,dk.
\]
第一式是 Fourier 反演公式,第二式是 Plancherel 公式.由于 Schwartz 空间在 \(L^2(\mathbb R^n)\) 中稠密,BLT 定理与本定理推出 Fourier 变换唯一地延拓为从 \(L^2(\mathbb R^n)\) 到自身的等距映射.
\end{theorem}

\begin{theorem}[Fourier 反演与 Plancherel 定理 II]\label{app:thm:plancherel-II}
Fourier 变换延拓为从 \(L^2(\mathbb R^n)\) 到 \(L^2(\mathbb R^n)\) 上的等距满射 \(\mathcal F\).它可按 \(L^2\) 范数拓扑中的极限计算:
\[
\mathcal F(\psi)(k)=(2\pi)^{-n/2}\lim_{A\to\infty}\int_{|x|\le A}e^{-ik\cdot x}\psi(x)\,dx,\tag{A.3}
\]
以及
\[
\mathcal F^{-1}f(x)=(2\pi)^{-n/2}\lim_{A\to\infty}\int_{|k|\le A}e^{ik\cdot x}f(k)\,dk.
\]
若 \(\psi\in L^1(\mathbb R^n)\cap L^2(\mathbb R^n)\),则由控制收敛定理,上式中的极限与定义 \ref{app:def:fourier} 中的 \(L^1\) Fourier 变换一致.
\end{theorem}

\begin{definition}[卷积]\label{app:def:convolution}
对两个可测函数 \(\phi,\psi\),定义卷积
\[
(\phi*\psi)(x)=\int_{\mathbb R^n}\phi(x-y)\psi(y)\,dy,
\]
其中假定对每个 \(x\) 积分绝对收敛.
\end{definition}

\begin{proposition}\label{app:prop:convolution-fourier}
若 \(\phi,\psi\in L^1(\mathbb R^n)\cap L^2(\mathbb R^n)\),则 \(\phi*\psi\) 有定义且属于 \(L^1(\mathbb R^n)\cap L^2(\mathbb R^n)\),并且
\[
(2\pi)^{-n/2}\mathcal F(\phi*\psi)=\mathcal F(\phi)\mathcal F(\psi).
\]
这个结果可由把 \(\phi*\psi\) 代入 Fourier 变换定义,写成 \(e^{-ik\cdot x}=e^{-ik\cdot y}e^{-ik\cdot(x-y)}\),再使用 Fubini 定理得到.
\end{proposition}

\begin{proposition}[Gaussian 积分]\label{app:prop:gaussian}
对任意 \(a>0\) 与 \(b\in\mathbb C\),有
\[
\frac1{\sqrt{2\pi}}\int_{-\infty}^{\infty}e^{-x^2/(2a)}e^{bx}\,dx
=\sqrt a\,e^{ab^2/2}.
\]
取 \(b=ik\) 可得到 Gaussian 函数 \(e^{-x^2/(2a)}\) 的 Fourier 变换;取 \(b=0\) 则可确定 Gaussian 概率密度的正确归一化.
\end{proposition}

\subsection{分布}

本节简要介绍分布理论——物理学家称之为“广义函数”——以及“分布意义下的导数”.基本思想是:通过考察函数与一类非常好的“测试函数”的积分来研究函数.例如,对局部可积函数 \(f\),考察
\[
\int_{\mathbb R^n}\chi(x)f(x)\,dx,\tag{A.4}
\]
其中 \(\chi\in C_c^\infty(\mathbb R^n)\) 是光滑紧支集函数.可以设想 \(\chi\) 非负、积分为 1,并且支集集中在某点 \(a\) 附近;此时上式是 \(f\) 在 \(a\) 处取值的一个近似,即物理学家所说的 \(f(a)\) 的“涂抹版本”.

\begin{proposition}\label{app:prop:test-function-uniqueness}
设 \(f_1,f_2\) 是 \(\mathbb R^n\) 上的局部可积函数.若对所有 \(\chi\in C_c^\infty(\mathbb R^n)\) 都有
\[
\int_{\mathbb R^n}\chi(x)f_1(x)\,dx
=\int_{\mathbb R^n}\chi(x)f_2(x)\,dx,
\]
则 \(f_1(x)=f_2(x)\) 几乎处处成立.
\end{proposition}

现在允许研究那些在点上未必有取值、但类似式 \(\ref{app:prop:test-function-uniqueness}\) 的表达式仍有意义的对象.从数学上说,我们把该表达式看作 \(C_c^\infty(\mathbb R^n)\) 上的线性泛函.

\begin{definition}[测试函数的收敛]\label{app:def:test-convergence}
若 \(\chi_m,\chi\in C_c^\infty(\mathbb R^n)\),并且满足:(1) 存在一个同时包含所有 \(\chi_m\) 支集的紧集 \(K\);(2) \(\chi_m\) 一致收敛到 \(\chi\);(3) \(\chi_m\) 的每阶导数都一致收敛到 \(\chi\) 的相应阶导数,则称 \(\chi_m\) 按测试函数意义收敛到 \(\chi\).
\end{definition}

\begin{definition}[分布]\label{app:def:distribution}
\(\mathbb R^n\) 上的一个分布是线性映射 \(T:C_c^\infty(\mathbb R^n)\to\mathbb C\),并满足如下连续性条件:若 \(\chi_m\) 按定义 \ref{app:def:test-convergence} 的意义收敛到 \(\chi\),则 \(T(\chi_m)\to T(\chi)\).
\end{definition}

这个连续性条件应当看作技术性要求:任何在整个 \(C_c^\infty(\mathbb R^n)\) 上定义良好且线性、并以相当自然的构造方式得到的泛函,都会满足该条件.

\begin{example}[Dirac \(\delta\)“函数”]\label{app:ex:dirac}
Dirac \(\delta\)“函数”是由
\[
\delta(\chi)=\chi(0)
\]
定义的分布.
\end{example}

\begin{definition}[由函数给出的分布]\label{app:def:distribution-function}
若 \(T\) 是分布,\(f\) 是局部可积函数,则称“\(T\) 等于 \(f\)”或“\(T\) 由 \(f\) 给出”,是指对所有 \(\chi\in C_c^\infty(\mathbb R^n)\),
\[
T(\chi)=\int_{\mathbb R^n}\chi(x)f(x)\,dx.
\]
\end{definition}

\begin{definition}[分布导数]\label{app:def:distribution-derivative}
若 \(T\) 是分布,定义分布 \(\partial T/\partial x_j\) 为
\[
\frac{\partial T}{\partial x_j}(\chi)
=-T\!\left(\frac{\partial\chi}{\partial x_j}\right).
\]
\end{definition}

容易验证,若 \(T\) 满足定义 \ref{app:def:distribution} 中的连续性条件,则 \(\partial T/\partial x_j\) 也满足该条件.此外,若 \(T\) 由连续可微函数给出,则 \(T\) 的分布导数与通常意义下的导数一致;这可以通过分部积分直接验证.对分布 \(T\),反复应用定义 \ref{app:def:distribution-derivative} 可以定义 \(\Delta T\),并得到
\[
(\Delta T)(\chi)=T(\Delta\chi).
\]

\begin{proposition}\label{app:prop:distribution-L2}
若 \(\phi,\psi\) 是 \(L^2\) 函数,则方程 \(\partial\psi/\partial x_j=\phi\) 在分布意义下成立,当且仅当对所有 \(\chi\in C_c^\infty(\mathbb R^n)\),
\[
-\left\langle\frac{\partial\chi}{\partial x_j},\psi\right\rangle
=\langle\chi,\phi\rangle.
\]
类似地,\(\Delta\psi=\phi\) 在分布意义下成立,当且仅当
\[
\langle\Delta\chi,\psi\rangle=\langle\chi,\phi\rangle.
\]
\end{proposition}

\begin{proposition}\label{app:prop:distribution-constant}
若 \(T\) 是 \(\mathbb R\) 上的分布,且 \(dT/dx\) 是零分布,则 \(T\) 是常数.也就是说,存在常数 \(c\),使得
\[
T(\chi)=\int_{-\infty}^{\infty}\chi(x)c\,dx.\tag{A.5}
\]
\end{proposition}

特别地,若 \(T\) 由局部可积函数 \(f\) 给出,且 \(T\) 的导数为零,那么存在常数 \(c\),使得对所有 \(\chi\in C_c^\infty(\mathbb R)\),
\[
\int_{-\infty}^{\infty}\chi(x)(f(x)-c)\,dx=0.
\]
由命题 \ref{app:prop:test-function-uniqueness},\(f(x)=c\) 几乎处处成立.这说明,即使 \(f\) 的导数只在弱意义(或分布意义)下为零,\(f\) 仍必须是常数.

\subsection{Banach 空间}\label{app:subsec:banach}

本节定义 Banach 空间并描述其若干初等性质.Banach 空间的一个重要例子是 \(L^p\) 空间.

\begin{definition}[范数]\label{app:def:norm}
设 \(V\) 是 \(\mathbb F\) 上的向量空间.范数是从 \(V\) 到 \(\mathbb R\) 的映射,记作 \(\psi\mapsto\|\psi\|\),满足:
\begin{enumerate}
  \item 对所有 \(\psi\in V\),\(\|\psi\|\ge0\),且等号成立当且仅当 \(\psi=0\);
  \item 对所有 \(\psi\in V,c\in\mathbb F\),\(\|c\psi\|=|c|\|\psi\|\);
  \item 对所有 \(\phi,\psi\in V\),\(\|\phi+\psi\|\le\|\phi\|+\|\psi\|\).
\end{enumerate}
若 \(\|\cdot\|\) 是 \(V\) 上的范数,则可定义距离 \(d(\phi,\psi)=\|\psi-\phi\|\).
\end{definition}

\begin{definition}[Banach 空间]\label{app:def:banach}
若赋予向量空间的距离空间关于上述距离完备,则称其为 Banach 空间.若 Banach 空间含有可数稠密子集,则称其可分.
\end{definition}

\begin{definition}[级数的收敛]\label{app:def:banach-series}
取值于赋范空间 \(V\) 的无穷级数 \(\sum_{n=1}^{\infty}\psi_n\) 称为收敛的,如果存在 \(L\in V\),使得令 \(S_N=\sum_{n=1}^N\psi_n\) 后有
\[
\lim_{N\to\infty}\|S_N-L\|=0.
\]
\end{definition}

\begin{proposition}[绝对收敛推出收敛]\label{app:prop:absolute-convergence}
若 \(V\) 是 Banach 空间,则绝对收敛推出在 \(V\) 中收敛.也就是说,若
\[
\sum_{n=1}^{\infty}\|\psi_n\|<\infty,
\]
则 \(\sum_{n=1}^{\infty}\psi_n\) 在 \(V\) 中收敛.
\end{proposition}

\begin{definition}[有界线性映射与算子范数]\label{app:def:bounded-map}
若 \(V_1,V_2\) 是赋范空间,线性映射 \(T:V_1\to V_2\) 称为有界的,如果
\[
\sup_{\psi\in V_1\setminus\{0\}}\frac{\|T\psi\|}{\|\psi\|}<\infty.\tag{A.6}
\]
若 \(T\) 有界,则式 \ref{app:def:bounded-map} 中的上确界称为算子范数,记作 \(\|T\|\).
\end{definition}

\begin{theorem}[有界线性变换定理]\label{app:thm:bounded-extension}
设 \(V_1\) 是赋范空间,\(V_2\) 是 Banach 空间,\(W\) 是 \(V_1\) 的稠密子空间,且 \(T:W\to V_2\) 是有界线性映射.则存在唯一的有界线性映射 \(\widetilde T:V_1\to V_2\),满足 \(\widetilde T|_W=T\).此外,\(\|\widetilde T\|=\|T\|\).
\end{theorem}

\begin{definition}[有界线性泛函与对偶空间]\label{app:def:dual}
设 \(V\) 是 \(\mathbb F\) 上的赋范空间.\(V\) 上的有界线性泛函是从 \(V\) 到 \(\mathbb F\) 的有界线性映射,其中 \(\mathbb F\) 使用绝对值给出的范数.所有有界线性泛函组成的空间,赋予由式 \ref{app:def:bounded-map} 给出的范数,称为 \(V\) 的对偶空间,记作 \(V^*\).
\end{definition}

\begin{theorem}\label{app:thm:dual-banach}
若 \(V\) 是赋范向量空间,则:
\begin{enumerate}
  \item 对偶空间 \(V^*\) 是 Banach 空间;
  \item 对所有 \(\psi\in V\),存在非零 \(\xi\in V^*\),使得 \(|\xi(\psi)|=\|\xi\|\|\psi\|\).特别地,若对所有 \(\xi\in V^*\) 都有 \(\xi(\psi)=0\),则 \(\psi=0\).
\end{enumerate}
\end{theorem}

\begin{theorem}[闭图定理]\label{app:thm:closed-graph}
设 \(V_1\) 是 Banach 空间,\(V_2\) 是赋范向量空间.对任意线性映射 \(T:V_1\to V_2\),令 \(\operatorname{Graph}(T)\) 表示 \(V_1\times V_2\) 中所有 \((\psi,T\psi)\)(\(\psi\in V_1\))组成的集合.若 \(T\) 的图是 \(V_1\times V_2\) 中的闭集,则 \(T\) 有界.
\end{theorem}

闭图定理的一个简单应用如下.设 \(V_1,V_2\) 是 Banach 空间,\(T:V_1\to V_2\) 是双射且有界的线性映射,则逆映射 \(T^{-1}:V_2\to V_1\) 自动有界.事实上,先检查 \(T\) 有界时其图是闭的;再注意到 \(T^{-1}\) 的图由 \(T\) 的图通过 \((\phi,\psi)\mapsto(\psi,\phi)\) 得到,因此也是闭的.闭图定理遂给出 \(T^{-1}\) 有界.

\begin{theorem}[一致有界原理]\label{app:thm:uniform-boundedness}
设 \(\{T_\alpha\}\) 是从 Banach 空间 \(V_1\) 到赋范空间 \(V_2\) 的任意一族有界线性映射.假定对每个 \(\psi\in V_1\),存在常数 \(C_\psi\),使得对所有 \(\alpha\),\(\|T_\alpha\psi\|\le C_\psi\).则存在常数 \(C\),使得对所有 \(\alpha\),\(\|T_\alpha\|\le C\).
\end{theorem}

这一定理的逆否命题表述为:若族 \(\{T_\alpha\}\) 无界,则必存在某个 \(\psi\in V_1\),使得 \(\{T_\alpha\psi\}\) 在 \(\psi\) 上无界.

\begin{corollary}\label{app:cor:bounded-set-dual}
设 \(V\) 是 Banach 空间,\(E\) 是 \(V\) 的非空子集.若对每个 \(\xi\in V^*\) 都存在常数 \(C_\xi\),使得对所有 \(\psi\in E\),\(|\xi(\psi)|\le C_\xi\),则 \(E\) 是有界集.
\end{corollary}

证明方法是把每个 \(\psi\in V\) 识别为 \(V^*\to\mathbb C\) 的线性映射 \(e_\psi\),即 \(e_\psi(\xi)=\xi(\psi)\).注意,由定理 \ref{app:thm:dual-banach} 的第 2 点,\(e_\psi\) 作为 \(V^{**}\) 中元素的范数等于 \(\psi\) 作为 \(V\) 中元素的范数.

\section{Hilbert 空间及其上的算子}

\subsection{内积空间与 Hilbert 空间}\label{app:subsec:hilbert}

现在把 \(\mathbb R^n\) 上通常的内积(或点积)推广到任意的 \(\mathbb R\) 或 \(\mathbb C\) 上向量空间.

\begin{definition}[内积]\label{app:def:inner-product}
域 \(\mathbb F\)(\(\mathbb F=\mathbb R\) 或 \(\mathbb C\))上向量空间 \(V\) 的内积是映射 \(\langle\cdot,\cdot\rangle:V\times V\to\mathbb F\),满足:
\begin{enumerate}
  \item 对所有 \(\phi,\psi\in V\),\(\langle\psi,\phi\rangle=\overline{\langle\phi,\psi\rangle}\);
  \item 对所有 \(\phi\in V\),\(\langle\phi,\phi\rangle\) 为实数且非负,并且 \(\langle\phi,\phi\rangle=0\) 当且仅当 \(\phi=0\);
  \item 对所有 \(\phi,\psi\in V,c\in\mathbb F\),\(\langle c\phi,\psi\rangle=\overline c\langle\phi,\psi\rangle\),\(\langle\phi,c\psi\rangle=c\langle\phi,\psi\rangle\);
  \item 对所有 \(\phi,\psi,\chi\in V\),\(\langle\phi+\psi,\chi\rangle=\langle\phi,\chi\rangle+\langle\psi,\chi\rangle\),且 \(\langle\phi,\psi+\chi\rangle=\langle\phi,\psi\rangle+\langle\phi,\chi\rangle\).
\end{enumerate}
\end{definition}

\begin{note}
我们遵循物理学约定,把复共轭放在内积的第一个因子上.
\end{note}

\begin{proposition}[Cauchy--Schwarz 不等式]\label{app:prop:cauchy-schwarz}
若 \(V\) 是内积空间,则对所有 \(\phi,\psi\in V\),
\[
|\langle\phi,\psi\rangle|^2\le\langle\phi,\phi\rangle\langle\psi,\psi\rangle.
\]
此外,若定义 \(\|\cdot\|:V\to\mathbb R\) 为
\[
\|\psi\|=\sqrt{\langle\psi,\psi\rangle},\tag{A.7}
\]
则 \(\|\cdot\|\) 是 \(V\) 上的范数.
\end{proposition}

\begin{definition}[Hilbert 空间]\label{app:def:hilbert}
Hilbert 空间是 \(\mathbb R\) 或 \(\mathbb C\) 上配有内积 \(\langle\cdot,\cdot\rangle\) 的向量空间 \(H\),并且 \(H\) 关于式 \ref{app:prop:cauchy-schwarz} 给出的范数完备.
\end{definition}

换言之,Hilbert 空间是范数来自内积的 Banach 空间.在本附录 A.4 中,\(H\) 可以表示任意 \(\mathbb R\) 或 \(\mathbb C\) 上的 Hilbert 空间;正文中则约定 \(H\) 表示可分复 Hilbert 空间.

\begin{definition}[Hilbert 空间直和]\label{app:def:hilbert-direct-sum}
设 \(\{H_j\}\) 是一列可分 Hilbert 空间.Hilbert 空间直和记为
\[
H:=\bigoplus_{j=1}^{\infty}H_j,\tag{A.8}
\]
它由所有序列 \(\psi=(\psi_1,\psi_2,\ldots)\) 构成,其中 \(\psi_n\in H_n\),且
\[
\|\psi\|^2:=\sum_{j=1}^{\infty}\|\psi_j\|_j^2<\infty.
\]
有限直和是所有除有限多个 \(\psi_j\) 外均为零的这类序列组成的集合.
\end{definition}

在直和上定义内积
\[
\langle\phi,\psi\rangle=\sum_{j=1}^{\infty}\langle\phi_j,\psi_j\rangle,\tag{A.9}
\]
其中 \(\phi,\psi\in H\).这个内积定义良好,并且 \(H\) 关于它完备,因此 \(H\) 是 Hilbert 空间.一个重要的 Hilbert 空间例子是测度空间 \((X,\mu)\) 上的 \(L^2(X,\mu)\).

\begin{definition}[\(L^2\) 空间的内积]\label{app:def:L2-inner}
若 \((X,\mu)\) 是测度空间,则在 \(L^2(X,\mu)\) 上定义
\[
\langle\phi,\psi\rangle=\int_X\overline{\phi(x)}\psi(x)\,d\mu(x).\tag{A.10}
\]
\end{definition}

测度论中的标准结果表明,式 \ref{app:def:L2-inner} 右端的积分对所有 \(\phi,\psi\in L^2(X,\mu)\) 都绝对收敛.于是容易验证 \(\langle\cdot,\cdot\rangle\) 确实是 \(L^2(X,\mu)\) 上的内积.另一个标准结果表明,\(L^2(X,\mu)\) 关于与该内积相联系的范数完备;因此 \(L^2(X,\mu)\) 是 Hilbert 空间.

\subsection{正交性}

Hilbert 空间比一般 Banach 空间更便于处理的一个原因,是 Hilbert 空间具有正交性概念.

\begin{definition}[正交]\label{app:def:orthogonal}
内积空间中的两个元素 \(\phi,\psi\) 称为正交的,如果 \(\langle\phi,\psi\rangle=0\).
\end{definition}

\begin{definition}[正交空间]\label{app:def:orthogonal-space}
若 \(V\) 是 \(H\) 的任意子空间,定义
\[
V^\perp=\{\phi\in H\mid\langle\phi,\psi\rangle=0\text{ 对所有 }\psi\in V\}.
\]
称 \(V^\perp\) 为 \(V\) 的正交空间.
\end{definition}

\begin{proposition}[正交分解]\label{app:prop:orthogonal-decomposition}
\begin{enumerate}
  \item 若 \(V\) 是 \(H\) 的闭子空间,则每个 \(\psi\in H\) 都能唯一分解为 \(\psi=\psi_1+\psi_2\),其中 \(\psi_1\in V,\psi_2\in V^\perp\);
  \item 若 \(V\) 是 \(H\) 的任意子空间,则 \((V^\perp)^\perp=\overline V\),其中 \(\overline V\) 是 \(V\) 的闭包.特别地,若 \(V\) 闭,则 \((V^\perp)^\perp=V\).
\end{enumerate}
若 \(V\) 闭,则称 \(V^\perp\) 为 \(V\) 的正交补.
\end{proposition}

\begin{definition}[正交归一集与正交归一基]\label{app:def:orthonormal-basis}
设 \(\{e_j\}\) 是 \(H\) 中的一族元素,其中 \(j\) 遍历任意指标集.若
\[
\langle e_j,e_k\rangle=
\begin{cases}0,&j\ne k,\\1,&j=k,\end{cases}
\]
则称其为正交归一集.若一组正交归一集的有限线性组合在 \(H\) 中稠密,则称它是 \(H\) 的正交归一基.
\end{definition}

若 \(H=L^2([-L,L])\),其中 \(L>0\),则函数
\[
\psi_n=\frac1{\sqrt{2L}}e^{2\pi inx/L},\qquad n\in\mathbb Z,\tag{A.11}
\]
构成 \(H\) 的一组正交归一基.

\begin{proposition}[正交归一展开]\label{app:prop:orthonormal-expansion}
设 \(\{e_j\}\) 是 \(H\) 的正交归一基,则每个 \(\psi\) 都能唯一表示为收敛和
\[
\psi=\sum_j a_je_j,\tag{A.12}
\]
其中 \(a_j=\langle e_j,\psi\rangle\).若 \(\psi\) 如式 \ref{app:prop:orthonormal-expansion} 所示,则
\[
\|\psi\|^2=\sum_j|a_j|^2.
\]
最后,若 \(\{a_j\}\) 是任意满足 \(\sum_j|a_j|^2<\infty\) 的数列,则存在唯一 \(\psi\in H\),使得对所有 \(j\),\(\langle e_j,\psi\rangle=a_j\).
\end{proposition}

当正交归一基取式 \ref{app:def:orthonormal-basis} 中的那一组时,由式 \ref{app:prop:orthonormal-expansion} 得到的级数称为 \(\psi\) 的 Fourier 级数.

\subsection{Riesz 定理与伴随算子}

记 \(B(H)\) 为从 \(H\) 到 \(H\) 的有界线性映射空间.不难证明,\(B(H)\) 关于算子范数构成 Banach 空间.

\begin{theorem}[Riesz 表示定理]\label{app:thm:riesz}
若 \(\xi:H\to\mathbb C\) 是有界线性泛函,则存在唯一的 \(\chi\in H\),使得对所有 \(\psi\in H\),
\[
\xi(\psi)=\langle\chi,\psi\rangle.
\]
此外,\(\xi\) 作为线性泛函的算子范数等于 \(\chi\) 作为 \(H\) 中元素的范数.
\end{theorem}

现在转向有界算子的伴随以及 Hilbert 空间上的二次型.

\begin{proposition}[伴随算子]\label{app:prop:adjoint}
对任意 \(A\in B(H)\),存在唯一的线性算子 \(A^*:H\to H\),称为 \(A\) 的伴随,使得
\[
\langle\phi,A\psi\rangle=\langle A^*\phi,\psi\rangle
\]
对所有 \(\phi,\psi\in H\) 成立.对 \(A,B\in B(H)\) 与 \(\alpha,\beta\in\mathbb C\),有
\[
(A^*)^*=A,\qquad (AB)^*=B^*A^*,\qquad
(\alpha A+\beta B)^*=\overline\alpha A^*+\overline\beta B^*,\qquad I^*=I.
\]
算子 \(A^*\) 有界,且 \(\|A^*\|=\|A\|\).
\end{proposition}

由于 \(A\) 有界,对每个固定的 \(\phi\in H\),映射 \(\psi\mapsto\langle\phi,A\psi\rangle\) 是有界线性泛函.由 Riesz 定理,存在唯一的 \(\chi\in H\),使 \(\langle\phi,A\psi\rangle=\langle\chi,\psi\rangle\).定义 \(A^*\phi=\chi\),即可得到 \(A^*\).不难检查,这一定义使 \(A^*\) 成为有界线性算子.

\begin{definition}[自伴与斜自伴]\label{app:def:self-adjoint-bounded}
算子 \(A\in B(H)\) 称为自伴的,如果 \(A^*=A\);称为斜自伴的,如果 \(A^*=-A\).
\end{definition}

\begin{definition}[酉算子]\label{app:def:unitary}
若 \(U\) 是 \(H\) 上的满射,并保持内积,即对所有 \(\phi,\psi\in H\),
\[
\langle U\phi,U\psi\rangle=\langle\phi,\psi\rangle,
\]
则称 \(U\) 为酉算子.
\end{definition}

若 \(U\) 是酉的,则它保持范数,\(\|U\psi\|=\|\psi\|\),因而有界且 \(\|U\|=1\).由极化恒等式(命题 \ref{app:prop:polarization}),保持范数也会推出保持内积.

\begin{proposition}[酉算子的判别]\label{app:prop:unitary-characterization}
有界算子 \(U\) 是酉算子,当且仅当 \(U^*=U^{-1}\),也就是当且仅当 \(UU^*=U^*U=I\).
\end{proposition}

\begin{proposition}[正交投影]\label{app:prop:orthogonal-projection}
对 \(H\) 的每个闭子空间 \(V\),存在唯一的有界算子 \(P\),使得 \(P\) 在 \(V\) 上等于 \(I\),在正交补 \(V^\perp\) 上等于 \(0\).该算子称为到 \(V\) 上的正交投影,并满足 \(P^2=P\)、\(P^*=P\).反过来,若 \(P\) 是 \(H\) 上满足 \(P^2=P\)、\(P^*=P\) 的有界算子,则 \(P\) 是到某个闭子空间 \(V\) 上的正交投影,其中 \(V=\operatorname{range}(P)\).
\end{proposition}

\subsection{二次型}

本节发展 Hilbert 空间上二次型的理论.由于通常只对内积本身作此处理,我们把所需结果的证明也包括在内.

\begin{definition}[半双线性型]\label{app:def:sesquilinear}
Hilbert 空间 \(H\) 上的半双线性型是映射 \(L:H\times H\to\mathbb C\),它关于第一个因子共轭线性,关于第二个因子线性.若存在常数 \(C\),使得对所有 \(\phi,\psi\in H\),
\[
|L(\phi,\psi)|\le C\|\phi\|\|\psi\|,
\]
则称 \(L\) 有界.
\end{definition}

\begin{proposition}[极化恒等式]\label{app:prop:polarization}
半双线性型 \(L\) 可以由其在对角线上的取值恢复:
\[
\begin{aligned}
L(\phi,\psi)={}&\frac12[L(\phi+\psi,\phi+\psi)-L(\phi,\phi)-L(\psi,\psi)]\\
&-\frac i2[L(\phi+i\psi,\phi+i\psi)-L(\phi,\phi)-L(i\psi,i\psi)].
\end{aligned}
\]
这个公式称为极化恒等式.这里不假定 \(L(\phi,\psi)\) 与 \(L(\psi,\phi)\) 之间存在任何关系.
\end{proposition}

\begin{proof}
直接计算即可.
\end{proof}

\begin{definition}[二次型]\label{app:def:quadratic-form}
Hilbert 空间 \(H\) 上的二次型是映射 \(Q:H\to\mathbb C\),满足:(1) 对所有 \(\psi\in H,\lambda\in\mathbb C\),\(Q(\lambda\psi)=|\lambda|^2Q(\psi)\);(2) 由
\[
\begin{aligned}
L(\phi,\psi)={}&\frac12[Q(\phi+\psi)-Q(\phi)-Q(\psi)]\\
&-\frac i2[Q(\phi+i\psi)-Q(\phi)-Q(i\psi)]
\end{aligned}
\]
定义的 \(L:H\times H\to\mathbb C\) 是半双线性型.若存在常数 \(C\),使 \(|Q(\phi)|\le C\|\phi\|^2\) 对所有 \(\phi\in H\) 成立,则称 \(Q\) 有界;满足此条件的最小 \(C\) 称为 \(Q\) 的范数.
\end{definition}

\begin{proposition}[二次型的性质]\label{app:prop:quadratic-properties}
设 \(Q\) 是 \(H\) 上的二次型,\(L\) 是相应的半双线性型,则:
\begin{enumerate}
  \item 对所有 \(\psi\in H\),\(Q(\psi)=L(\psi,\psi)\);
  \item 若 \(Q\) 有界,则 \(L\) 有界;
  \item 若对所有 \(\psi\in H\),\(Q(\psi)\in\mathbb R\),则 \(L\) 共轭对称,即 \(L(\phi,\psi)=\overline{L(\psi,\phi)}\).
\end{enumerate}
\end{proposition}

\begin{proof}
第 1 点可在 \(L(\phi,\psi)\) 的表达式中取 \(\phi=\psi\),再使用 \(Q(\lambda\psi)=|\lambda|^2Q(\psi)\) 验证.为证明第 2 点,设 \(|Q(\psi)|\le C\|\psi\|^2\).若 \(\|\phi\|=\|\psi\|=1\),则 \(\|\phi+\psi\|\) 与 \(\|\phi+i\psi\|\) 至多为 2,因此 \(|L(\phi,\psi)|\le6C\).

对任意 \(\phi,\psi\),写成 \(\phi=\|\phi\|\widetilde\phi,\psi=\|\psi\|\widetilde\psi\),其中带波浪号的向量为单位向量.由半双线性性即得
\[
|L(\phi,\psi)|\le6C\|\phi\|\|\psi\|.
\]
最后,若 \(Q\) 处处取实值,定义
\[
M(\phi,\psi)=\frac12[Q(\phi+\psi)-Q(\phi)-Q(\psi)]=\operatorname{Re}L(\phi,\psi).
\]
则 \(M\) 是实双线性的并且对称,且 \(M(i\phi,i\psi)=M(\phi,\psi)\).由此 \(M(\phi,i\psi)=-M(\psi,i\phi)\),从而
\[
L(\phi,\psi)=M(\phi,\psi)-iM(\phi,i\psi)
=\overline{L(\psi,\phi)}.
\]
\end{proof}

\begin{example}[由有界算子构造二次型]\label{app:ex:operator-quadratic}
若 \(A\) 是 \(H\) 上的有界算子,则可定义有界二次型
\[
Q_A(\psi)=\langle\psi,A\psi\rangle,\qquad\psi\in H.
\]
相应的半双线性型为
\[
L_A(\phi,\psi)=\langle\phi,A\psi\rangle.
\]
\end{example}

\begin{proposition}[二次型的算子表示]\label{app:prop:quadratic-operator}
若 \(Q\) 是 \(H\) 上的有界二次型,则存在唯一的 \(A\in B(H)\),使得对所有 \(\psi\in H\),\(Q(\psi)=\langle\psi,A\psi\rangle\).若 \(Q(\psi)\) 对所有 \(\psi\) 都是实数,则 \(A\) 自伴.
\end{proposition}

\begin{proof}
由命题 \ref{app:prop:quadratic-properties},\(L\) 有界,即存在 \(C\) 使 \(|L(\phi,\psi)|\le C\|\phi\|\|\psi\|\).固定 \(\phi\),则 \(\psi\mapsto L(\phi,\psi)\) 是范数不超过 \(C\|\phi\|\) 的有界线性泛函.由 Riesz 定理,存在唯一的 \(\chi\in H\),满足 \(\|\chi\|\le C\|\phi\|\) 且 \(L(\phi,\psi)=\langle\chi,\psi\rangle\).定义 \(B\phi=\chi\).直接计算表明 \(B\) 线性,而不等式表明 \(B\) 有界.令 \(A=B^*\),便得到所需算子.唯一性来自如下观察:若对所有 \(\phi,\psi\) 都有 \(\langle\phi,A\psi\rangle=0\),则 \(A\) 是零算子.

若 \(Q\) 处处取实值,则由命题 \ref{app:prop:quadratic-properties} 的第 3 点,\(L\) 共轭对称.因此
\[
\langle\phi,A\psi\rangle=L(\phi,\psi)=\overline{L(\psi,\phi)}
=\overline{\langle\psi,A\phi\rangle}=\langle A\phi,\psi\rangle,
\]
故 \(A\) 自伴.
\end{proof}

\subsection{Hilbert 空间的张量积}\label{app:subsec:hilbert-tensor}

回忆附录 \ref{app:def:tensor} 中两个向量空间的张量积概念.

\begin{proposition}[内积的张量积]\label{app:prop:inner-tensor}
设 \(V_1,V_2\) 是内积空间,内积分别为 \(\langle\cdot,\cdot\rangle_1\) 与 \(\langle\cdot,\cdot\rangle_2\).则在 \(V_1\otimes V_2\) 上存在唯一的内积 \(\langle\cdot,\cdot\rangle\),满足
\[
\langle u_1\otimes v_1,u_2\otimes v_2\rangle
=\langle u_1,u_2\rangle_1\langle v_1,v_2\rangle_2
\]
对所有 \(u_1,u_2\in V_1,v_1,v_2\in V_2\) 成立.
\end{proposition}

若 \(H_1,H_2\) 是 Hilbert 空间,可以用命题 \ref{app:prop:inner-tensor} 在代数张量积 \(H_1\otimes H_2\) 上定义内积.然而,若 \(H_1,H_2\) 都是无限维的,\(H_1\otimes H_2\) 关于这个内积通常并不完备.我们可以关于该内积完成 \(H_1\otimes H_2\),从而得到新的 Hilbert 空间.

\begin{definition}[Hilbert 空间张量积]\label{app:def:hilbert-tensor}
若 \(H_1,H_2\) 是 Hilbert 空间,则它们的 Hilbert 空间张量积,记作 \(H_1\widehat\otimes H_2\),是关于命题 \ref{app:prop:inner-tensor} 中内积完成 \(H_1\otimes H_2\) 所得到的 Hilbert 空间.
\end{definition}

\begin{proposition}[张量积中的正交归一基]\label{app:prop:hilbert-tensor-basis}
若 \(H_1,H_2\) 是 Hilbert 空间,分别具有正交归一基 \(\{e_j\}\)、\(\{f_k\}\),则 \(\{e_j\otimes f_k\}\) 是 \(H_1\widehat\otimes H_2\) 的一组正交归一基.
\end{proposition}

\begin{proposition}[Hilbert 空间上的算子张量积]\label{app:prop:hilbert-operator-tensor}
若 \(A\) 是 \(H_1\) 上的有界算子,\(B\) 是 \(H_2\) 上的有界算子,则存在 \(H_1\widehat\otimes H_2\) 上唯一的有界算子 \(A\otimes B\),满足
\[
(A\otimes B)(\phi\otimes\psi)=(A\phi)\otimes(B\psi)
\]
对所有 \(\phi\in H_1,\psi\in H_2\) 成立.此外,\(\|A\otimes B\|=\|A\|\|B\|\).
\end{proposition}

为说明有界性,先写 \(A\otimes B=(A\otimes I)(I\otimes B)\).取 \(H_2\) 的任意正交归一基 \(\{f_j\}\),可以把 \(H_1\widehat\otimes H_2\) 分解为形如 \(H_1\otimes f_j\) 的子空间的 Hilbert 空间直和.算子 \(A\otimes I\) 在这个分解下是每个对角块均为 \(A\) 的分块对角算子,因此 \(\|A\otimes I\|=\|A\|\).类似地,\(\|I\otimes B\|=\|B\|\),所以
\[
\|A\otimes B\|\le\|A\otimes I\|\|I\otimes B\|=\|A\|\|B\|.
\]
另一方面,取单位向量列 \(\phi_n\in H_1,\psi_n\in H_2\),使 \(\|A\phi_n\|\to\|A\|\)、\(\|B\psi_n\|\to\|B\|\),便得到反向不等式,因而等号成立.
